% GENERATED FILE. Edit canonical lessons, then run npm run generate:books.
% canonical-source-hash: fnv1a64:b3d925018005e2eb
% canonical-lessons: HI-C381-kas, HI-C381-afsos, HI-C381-satana, HI-C381-gavana, HI-C381-bhala

\chapter{If Only}
\label{ch:hi-if-only}
\hlchaptermodality{381}

\begin{chapteropening}
\textbf{By the end of this chapter:} \emph{I can say what I wish had happened.}

Complete the last lesson of chapter 381: I can say what I wish had happened.
\end{chapteropening}

\section[\texorpdfstring{\hi{काश} (kāś)}{काश (kāś)}]{\hi{काश} (kāś) — if only, I wish}

\label{lesson:HI-C381-kas}



\begin{quote}
Before the new one: say the Hindi for to get stuck, then the Hindi for to pull out of trouble.
\end{quote}

\subsection*{You'll want to know: \hi{काश}}
\textbf{\hi{काश}} — \emph{kāś} — \textquotedblleft{}if only, I wish\textquotedblright{}.

\begin{quote}
\hi{काश} \hi{मैं} \hi{वहाँ} \hi{होता}, \hi{तो} \hi{तुम्हारी} \hi{मदद} \hi{करता}!
\end{quote}

\begin{quote}
\emph{kāś maiṁ vahā̃ hotā, to tumhārī madad kartā!}
\end{quote}

\begin{quote}
If only I had been there, I would have helped you!
\end{quote}

\begin{cousinweb}
\textbf{\hi{काश}} (\emph{kāś}) is usually traced to Persian \emph{kāsh}, \textquotedblleft{}would that\textquotedblright{}.
\end{cousinweb}

\begin{grammarlens}[title={What you've built}]
One more word for saying what happens if.
\end{grammarlens}

\subsection*{Your turn}
\begin{itemize}
\raggedright
  \item \emph{Say it:} \emph{kāś}
  \item \emph{Say it:} \emph{kāś}, once more
  \item \emph{Say it:} \emph{kāś maiṁ vahā̃ hotā, to tumhārī madad kartā!} — If only I had been there, I would have helped you!
  \item \emph{Recall:} say the Hindi for to get stuck, then the Hindi for to pull out of trouble, then say \emph{kāś} again
\end{itemize}

\begin{tcolorbox}[breakable,colback=teal!4,colframe=teal!35!black,title={Before you move on}]
What does \hi{काश} mean? (\textquotedblleft{}If only, I wish\textquotedblright{}.) Say it once more.
\end{tcolorbox}
\section[\texorpdfstring{\hi{अफ़सोस} (afsos)}{अफ़सोस (afsos)}]{\hi{अफ़सोस} (afsos) — regret, a pity}

\label{lesson:HI-C381-afsos}



\begin{quote}
Before the new one: say the Hindi for to pull out of trouble, then the Hindi for if only.
\end{quote}

\subsection*{You'll want to know: \hi{अफ़सोस}}
\textbf{\hi{अफ़सोस}} — \emph{afsos} — \textquotedblleft{}regret, a pity\textquotedblright{}.

\begin{quote}
\hi{अफ़सोस} \hi{है} \hi{कि} \hi{अगर} \hi{मैं} \hi{जल्दी} \hi{आता} \hi{तो} \hi{मीरा} \hi{से} \hi{मिलता।}
\end{quote}

\begin{quote}
\emph{afsos hai ki agar maiṁ jaldī ātā to Mīrā se miltā.}
\end{quote}

\begin{quote}
It's a pity: if I had come earlier, I would have met Mira.
\end{quote}

Masculine.

\begin{cousinweb}
\textbf{\hi{अफ़सोस}} (\emph{afsos}) is usually traced to Persian \emph{afsōs}, \textquotedblleft{}regret, alas\textquotedblright{}.
\end{cousinweb}

\begin{grammarlens}[title={What you've built}]
One more word for saying what happens if.
\end{grammarlens}

\subsection*{Your turn}
\begin{itemize}
\raggedright
  \item \emph{Say it:} \emph{afsos}
  \item \emph{Say it:} \emph{afsos}, once more
  \item \emph{Say it:} \emph{afsos hai ki agar maiṁ jaldī ātā to Mīrā se miltā.} — It's a pity: if I had come earlier, I would have met Mira.
  \item \emph{Recall:} say the Hindi for to pull out of trouble, then the Hindi for if only, then say \emph{afsos} again
\end{itemize}

\begin{tcolorbox}[breakable,colback=teal!4,colframe=teal!35!black,title={Before you move on}]
What does \hi{अफ़सोस} mean? (\textquotedblleft{}Regret, a pity\textquotedblright{}.) Say it once more.
\end{tcolorbox}
\section[\texorpdfstring{\hi{सताना} (satānā)}{सताना (satānā)}]{\hi{सताना} (satānā) — to torment, to harass}

\label{lesson:HI-C381-satana}



\begin{quote}
Before the new one: say the Hindi for if only, then the Hindi for a pity.
\end{quote}

\subsection*{You'll want to know: \hi{सताना}}
\textbf{\hi{सताना}} — \emph{satānā} — \textquotedblleft{}to torment, to harass\textquotedblright{}.

\begin{quote}
\hi{अगर} \hi{तुम} \hi{उसे} \hi{नहीं} \hi{सताते} \hi{तो} \hi{वह} \hi{नहीं} \hi{रोता।}
\end{quote}

\begin{quote}
\emph{agar tum use nahīṃ satāte to vah nahīṃ rotā.}
\end{quote}

\begin{quote}
If you hadn't tormented him, he wouldn't have cried.
\end{quote}

Past, with \hi{ने} (he or she …): \hi{उसने} \hi{छोटे} \hi{भाई} \hi{को} \hi{सताया}, \emph{usne choṭe bhāī ko satāyā}.

\begin{cousinweb}
\textbf{\hi{सताना}} (\emph{satānā}) is a tadbhava, inherited from Sanskrit \emph{saṃtāpayati}, \textquotedblleft{}inflames; torments\textquotedblright{} (CDIAL 12886).
\end{cousinweb}

\begin{grammarlens}[title={What you've built}]
One more word for saying what happens if.
\end{grammarlens}

\subsection*{Your turn}
\begin{itemize}
\raggedright
  \item \emph{Say it:} \emph{satānā}
  \item \emph{Say it:} \emph{satānā}, once more
  \item \emph{Say it:} \emph{agar tum use nahīṃ satāte to vah nahīṃ rotā.} — If you hadn't tormented him, he wouldn't have cried.
  \item \emph{Recall:} say the Hindi for if only, then the Hindi for a pity, then say \emph{satānā} again
\end{itemize}

\begin{tcolorbox}[breakable,colback=teal!4,colframe=teal!35!black,title={Before you move on}]
What does \hi{सताना} mean? (\textquotedblleft{}To torment, to harass\textquotedblright{}.) Say it once more.
\end{tcolorbox}
\section[\texorpdfstring{\hi{गँवाना} (gãvānā)}{गँवाना (gãvānā)}]{\hi{गँवाना} (gãvānā) — to waste, to throw away (time, a chance)}

\label{lesson:HI-C381-gavana}



\begin{quote}
Before the new one: say the Hindi for a pity, then the Hindi for to torment.
\end{quote}

\subsection*{You'll want to know: \hi{गँवाना}}
\textbf{\hi{गँवाना}} — \emph{gãvānā} — \textquotedblleft{}to waste, to throw away (time, a chance)\textquotedblright{}.

\begin{quote}
\hi{अगर} \hi{तुम} \hi{खेल} \hi{में} \hi{समय} \hi{गँवाओगे} \hi{तो} \hi{काम} \hi{पूरा} \hi{नहीं} \hi{होगा।}
\end{quote}

\begin{quote}
\emph{agar tum khel meṁ samay gãvāoge to kām pūrā nahīṃ hogā.}
\end{quote}

\begin{quote}
If you waste time playing, the work won't get done.
\end{quote}

Past, with \hi{ने} (he or she …): \hi{उसने} \hi{पूरा} \hi{दिन} \hi{गँवाया}, \emph{usne pūrā din gãvāyā}.

\begin{cousinweb}
\textbf{\hi{गँवाना}} (\emph{gãvānā}) is a tadbhava, inherited from Sanskrit \emph{gamayati}, \textquotedblleft{}makes go\textquotedblright{} (CDIAL 4028).
\end{cousinweb}

\begin{grammarlens}[title={What you've built}]
One more word for saying what happens if.
\end{grammarlens}

\subsection*{Your turn}
\begin{itemize}
\raggedright
  \item \emph{Say it:} \emph{gãvānā}
  \item \emph{Say it:} \emph{gãvānā}, once more
  \item \emph{Say it:} \emph{agar tum khel meṁ samay gãvāoge to kām pūrā nahīṃ hogā.} — If you waste time playing, the work won't get done.
  \item \emph{Recall:} say the Hindi for a pity, then the Hindi for to torment, then say \emph{gãvānā} again
\end{itemize}

\begin{tcolorbox}[breakable,colback=teal!4,colframe=teal!35!black,title={Before you move on}]
What does \hi{गँवाना} mean? (\textquotedblleft{}To waste, to throw away (time, a chance)\textquotedblright{}.) Say it once more.
\end{tcolorbox}
\section[\texorpdfstring{\hi{भला} (bhalā)}{भला (bhalā)}]{\hi{भला} (bhalā) — good, decent}

\label{lesson:HI-C381-bhala}



\begin{quote}
Before the new one: say the Hindi for to torment, then the Hindi for to waste.
\end{quote}

\subsection*{You'll want to know: \hi{भला}}
\textbf{\hi{भला}} — \emph{bhalā} — \textquotedblleft{}good, decent\textquotedblright{}.

\begin{quote}
\hi{अगर} \hi{तुम} \hi{भले} \hi{आदमी} \hi{होते} \hi{तो} \hi{सच} \hi{बोलते।}
\end{quote}

\begin{quote}
\emph{agar tum bhale ādmī hote to sac bolte.}
\end{quote}

\begin{quote}
If you were a decent man, you would tell the truth.
\end{quote}

\begin{cousinweb}
\textbf{\hi{भला}} (\emph{bhalā}) is a tadbhava, inherited from Sanskrit \emph{bhalla}, \textquotedblleft{}auspicious\textquotedblright{} (CDIAL 9408).
\end{cousinweb}

\begin{grammarlens}[title={What you've built}]
One more word for saying what happens if.
\end{grammarlens}

\subsection*{Your turn}
\begin{itemize}
\raggedright
  \item \emph{Say it:} \emph{bhalā}
  \item \emph{Say it:} \emph{bhalā}, once more
  \item \emph{Say it:} \emph{agar tum bhale ādmī hote to sac bolte.} — If you were a decent man, you would tell the truth.
  \item \emph{Recall:} say the Hindi for to torment, then the Hindi for to waste, then say \emph{bhalā} again
\end{itemize}

\begin{tcolorbox}[breakable,colback=teal!4,colframe=teal!35!black,title={Before you move on}]
What does \hi{भला} mean? (\textquotedblleft{}Good, decent\textquotedblright{}.) Say it once more.
\end{tcolorbox}
